\chapter{Motivation}\label{ch:motivation}

\section{Why spot spraying needs a weed-versus-crop decision}
Weeds compete with the crop for light, water and nutrients, and herbicide is the cheapest way to
control them. It is usually applied uniformly across a field, although weeds cover only part of it.
Site-specific weed management (SSWM) instead treats only the places where weeds grow; reviews of the
field report that the saving in herbicide depends strongly on how patchy the weeds
are~\cite{gerhards2022sswm}. The finest form of SSWM is \emph{spot spraying}: a camera looks at the
ground ahead of a boom, and each nozzle opens only over a weed. Low-cost prototypes of such sprayers
already exist~\cite{partel2019sprayer}.

The hard part of spot spraying is not the nozzle but the decision. Slaughter et al.\ identified
reliable weed identification as the main obstacle to robotic weed control more than fifteen years
ago~\cite{slaughter2008review}, and it still is. The two possible mistakes are not equal: a missed
weed survives until the next pass, but a false ``weed'' on a crop plant damages the crop the farmer is
trying to protect. A spot sprayer is therefore only useful if its vision system keeps the crop off
the spray map.

\section{Why zero-shot vision}
Today this separation is done by supervised networks trained on thousands of labelled images of the
same crop under similar conditions~\cite{milioto2018realtime,lottes2018fcn,hasan2021survey}. Labelling
plants pixel by pixel is slow and expensive, and a network trained on one field often degrades on
another~\cite{weyler2023domain}. For a small Indian farm growing several crops in a season, collecting
a labelled dataset for each crop and field is not realistic.

Foundation models offer a different route. Grounding DINO~\cite{liu2024groundingdino} finds objects
named in a free-text prompt, and SAM~2~\cite{ravi2025sam2} outlines any object it is pointed at.
Chained together, as in Grounded SAM~\cite{ren2024groundedsam}, they promise to segment ``the weeds''
in a photograph without a single training label. Whether this works for crop and weed, which are both
small green plants that often look alike, is an open question: a 2025 study by Chong et
al.~\cite{chong2025zeroshot} argues that text prompts are ill-defined for exactly this task. AgriAgent
measures the question on real field photographs and asks what a transparent safety layer can recover.

\section{Problem statement}
Given an RGB field photograph and the name of the crop, produce a spray map that covers the weeds and
leaves the crop untouched, using pretrained models only (no training on our data), on a laptop GPU,
with every spraying decision explainable and testable.

\section{Scope}
AgriAgent is a research prototype. It produces a map on screen and a GeoJSON file; it does not drive a
sprayer, a drone or a tractor, and its GeoJSON coordinates are placeholders, not GPS. Separating weed
from crop for spot spraying is the only application studied. No model is trained or fine-tuned.
Disease detection, yield counting, canopy cover and nutrient stress, listed in our original proposal,
are future scope (Chapter~\ref{ch:future}).

\section{Contributions}
\begin{enumerate}
\item An honest measurement of the Grounded SAM pipeline for crop-versus-weed spraying on two crops,
  with labels copied from the official dataset files and a held-out set evaluated once
  (Chapter~\ref{ch:results}).
\item A deterministic safety gate and a four-rule plant classifier that move the held-out result from
  65.5\% crop pixels sprayed and 1.8\% weed recall to 8.8\% and 77.9\%, with no training and no model
  calls (Section~\ref{sec:holdout}).
\item A bounded escalation agent whose verdicts are combined by the measured accuracy of each source,
  and a negative but useful result: it is safe and precise but adds no weed recall
  (Section~\ref{sec:agent-results}).
\item A testing record: 270 automated tests, two system test cycles by a team member on a fresh
  Windows installation, a verification run, a manual test cycle and 30 findings (F1--F30) traced to
  fixes or stated limitations (Section~\ref{sec:testing}). The final build is released as
  \code{v1.0.0}.
\end{enumerate}

\section{Organisation of the report}
Chapter~\ref{ch:survey} surveys related work and Chapter~\ref{ch:gaps} states the gaps it leaves.
Chapter~\ref{ch:objectives} lists our objectives and research questions. Chapter~\ref{ch:proposed}
describes the proposed system, the datasets and the implementation. Chapter~\ref{ch:results} gives the
mathematical model, the evaluation method, the results with their analysis, and the testing.
Chapter~\ref{ch:conclusion} concludes and Chapter~\ref{ch:future} describes future scope.
Appendix~\ref{app:run} explains how to run the system; Appendix~\ref{app:contrib} records who built
which part, counted from the repository history.
